%----------------------------------------------------------------------------------------
% ДҮГНЭЛТ
%----------------------------------------------------------------------------------------
\chapter{Дүгнэлт}

\section{Судалгааны үр дүнгийн хураангуй}

Энэхүү дипломын ажилд хяналтын камерын бичлэгээс зодооны зөрчлийг автоматаар илрүүлж, Монгол хэлээр бүтэцтэй тайлбар гаргадаг багш--сурагч загварын мэдлэг дамжуулалт систем боловсруулсан. Судалгааны явцад хийгдсэн гол ажлуудыг хураангуйлбал:

\begin{itemize}
    \item \textbf{Сурагч загварын нарийн тохируулга:} Qwen2.5-VL-7B загварыг LoRA аргаар нарийвчлан тохируулж, зодооны хоёртын ангиллын даалгаварт туршсан.

    \item \textbf{Сургалтын хяналт:} WandB болон Rich terminal хяналтын самбараар сургалтын явцыг бодит цагаар хянаж, 500 алхам бүрт хадгалалтын цэг үүсгэсэн. 5-р үед сургалтын алдагдал $\sim$0.38*, баталгаажуулалтын алдагдал $\sim$0.64*, токений нарийвчлал $\sim$93\%* хүрсэн.

    \item \textbf{Туршилтын үр дүн:} Val өгөгдөл дээрх хоёртын ангиллын загварын үнэлгээнд stillmeta/fight-video-qlora1 загвар Accuracy=0.9179, F1=0.9238, Precision=0.9238, Recall=0.9238 үр дүн үзүүлсэн. Confusion matrix тэнцвэртэй гарч, Fight ангилалд 97 зөв/8 алдаа, No Fight ангилалд 82 зөв/8 алдаа бүртгэгдсэн. Энэ нь Qwen2.5-VL-7B BASE загварын 54.36\% accuracy-аас 91.79\% болж өссөн буюу нарийвчилсан тохируулгын бодит нөлөөг харуулсан.

\end{itemize}

\section{Дипломын зорилгуудын биелэлт}

\begin{table}[htbp]
\centering
\caption{Дипломын зорилгуудын биелэлтийн хүснэгт}
\label{tab:objectives}
\begin{tabular}{p{7cm}p{6cm}}
\hline
\textbf{Зорилго} & \textbf{Биелэлт} \\
\hline
Өгөгдлийн чанараар $\geq$10\% сайжруулалт & Өгөгдлийн урьдчилсан боловсруулалт хийсэн \checkmark \\
Хөнгөн тохируулгатай загвар боловсруулах & LoRA adapter ашигласан \checkmark \\
3+ алгоритм харьцуулалт                   & 3 загвар харьцуулсан \checkmark \\
Гиперпараметрийн оновчлол                 & Сургалтын тохиргоог туршилтаар сонгосон \checkmark \\
Transfer learning, нөөц $\geq$30\% хэмнэх & Сурагч загвар ашигласан \checkmark \\
Синтетик өгөгдөл үүсгэх                   & Ашиглаагүй \\
Загварын хяналт                           & WandB + хадгалалтын цэг \checkmark \\
Алдааны шинжилгээ                         & FP/FN шинжилгээ \checkmark \\
\end{tabular}
\end{table}

%-------------------------------------------------------------------------------
\section{Онолын хувь нэмэр}
%-------------------------------------------------------------------------------

Энэхүү судалгааны онолын хувь нэмрийг дараах байдлаар тодорхойлж болно:

\textbf{1. Далд мэдлэг дамжуулалт аргыг видео тандалтад хэрэгжүүлсэн.} Уламжлалт мэдлэг дамжуулалт арга нь soft logit-ийг шууд дамжуулдаг бол энэхүү судалгаа нь багш загварын бэлэн гаралтыг сурагч загварын сургалтын зорилт болгон ашигладаг offline KD аргыг видео тандалтын хүрээнд хэрэгжүүлсэн анхны судалгааны ажлуудын нэг болж байна. Энэ арга нь сургалтын явцад багш загварыг дахин ажиллуулах шаардлагагүй тул нөөцийн хэрэглээг эрс бууруулдаг.

\textbf{2. Автоматжуулсан Монгол хэлний тэмдэглэгээний арга зүй.} Хүний гараар шошголохын оронд багш VLM ашиглан Монгол хэлний Q\&A өгөгдлийн багц автоматаар үүсгэх арга зүйг санал болгосон. Энэ нь бага нөөцтэй хэлний NLP өгөгдлийн багц бэлтгэхэд өргөн хэрэглэгдэх арга зүйн загвар болох боломжтой.

\textbf{3. LoRA + QLoRA-ийн видео ойлголтод хэрэглэх.} Зодооны илрүүлэлт болон Монгол хэлний тайлбар гаргах хоёр зорилтыг нэгдмэл байдлаар шийдэхэд LoRA тохируулагчийн параметрийн тоо болон rank-ийн оновчтой утгыг туршилтаар тодорхойлсон нь цаашдын судалгаанд суурь болно.

\textbf{4. 4 түвшний зөрчлийн ангилал.} Уламжлалт хүчирхийлэлтэй/хүчирхийлэлгүй хоёртын ангилалтыг байхгүй, бага, дунд, өндөр гэсэн 4 түвшинд өргөтгөсөн нь аюулгүй байдлын байгууллагуудын яаралтай байдлын тодорхойлолтод практикт ашиглах боломжтой шинэ стандарт болж байна.

%-------------------------------------------------------------------------------
\section{Практик ач холбогдол}
%-------------------------------------------------------------------------------

Энэхүү судалгааны үр дүнгийн практик ач холбогдлыг дараах хэд хэдэн чиглэлд авч үзэж болно.

\textbf{Аюулгүй байдлын байгууллагуудад:} Нарийн тохируулсан загвар нь хяналтын камерын бичлэгийг боловсруулах боломжтой. Монгол хэлний бүтэцтэй тайлбар нь цагдаагийн ажилтанд хэрэг шалгах явцыг хурдасгаж, нотлох баримт бэлдэхэд шууд дэмжлэг үзүүлнэ.

\textbf{Тооцооллын нөөцийн хэмнэлт:} багш загварын 62.2 GB VRAM-ийн оронд сурагч загвар зөвхөн 6.5 GB ашигладаг. Энэ нь системийг байгууллагын түвшний GPU байхгүй орчинд, жишээлбэл NVIDIA RTX 3090 (24 GB) зэрэг нийтлэг GPU дээр ажиллуулах боломжийг олгодог.

\textbf{Арга зүйн давтан хэрэглэх боломж:} Энэ ажлын багш--сурагч загварын урсгал болон автомат тэмдэглэгээний арга зүй нь зодооны илрүүлэлтийн хүрээнд хязгаарлагдахгүйгээр аливаа видео ойлголтын даалгаварт (гал, ослын илрүүлэлт г.м.) хэрэглэх боломжтой.

%-------------------------------------------------------------------------------
\section{Нийгмийн нөлөөлөл}
%-------------------------------------------------------------------------------

Автоматчилсан зодооны илрүүлэлтийн системийн нийгмийн нөлөөлөл нь хоёрдмол шинж чанартай. Нэг талаас, систем нь:

\begin{itemize}
    \item Цагдаагийн хариу арга хэмжээний хугацааг багасгаж, иргэдийн аюулгүй байдлыг дэмжинэ.
    \item Хүний оператор дутмаг байх нөхцөлд хяналтын үр ашгийг нэмэгдүүлнэ.
    \item Нотлох баримт бэлдэх явцыг автоматжуулж, хэрэг шалгах ажлын ачааллыг хөнгөвчилнэ.
    \item Боловсролын байгууллага, нийтийн тээвэр зэрэг олон орчинд аюулгүй байдлын хяналтыг сайжруулна.
\end{itemize}

Нөгөө талаас, хяналтын камерын системийг өргөжүүлэх нь нийгмийн хяналтын (mass surveillance) эрсдэлтэй. Энэ эрсдэлийг бууруулахын тулд систем хэрэгжүүлэхэд дараах зарчмуудыг заавал баримтлах хэрэгтэй: хандалтын хяналт, аудитын мөр, нүүрийг анонимчлах, хуулийн зохицуулалтын хүрээнд ажиллах.

%-------------------------------------------------------------------------------
\section{Судалгааны хязгаарлалт ба хүчин чадал}
%-------------------------------------------------------------------------------

\subsection{Загварын хязгаарлалтууд}

\begin{enumerate}
    \item \textbf{Монгол хэлний өгөгдөл цөөн:} Qwen загварын хувилбаруудад Монгол хэлний урьдчилсан сургалтын өгөгдөл цөөн тул тайлбарын хэлний чанар англи хэлтэй харьцуулахад доогуур байдаг. Нэмэлт Монгол хэлний хяналтын камерын өгөгдлийн багц цуглуулах нь энэ хязгаарлалтыг хэсэгчлэн шийднэ.

    \item \textbf{Нэгэн зэрэг олон камер боловсруулах чадваргүй:} Одоогийн систем нэг цагийн урсгалд нэг загвар ажиллуулдаг. Олон камерын хувьд параллель процессорын дизайн шаардлагатай.

    \item \textbf{Тайлбар тайланд суурилсан оролт:} Багш загварын дүгнэлт гаргалт нь видео кадрын оронд тайлбар тайланд суурилсан тул видеоны баялаг мэдээллийн зарим хэсэг алдагдана. Шууд видео оролт бүхий багш загварын дүгнэлт гаргалт нь илүү оновчтой Q\&A өгөгдлийн багц үүсгэх боломжтой.

\end{enumerate}

\subsection{Судалгааны хүчин чадал}

\begin{itemize}
    \item \textbf{Өгөгдлийн багцын олон янз байдал:} Нийт 1,951 клипийг 80/10/10 харьцаагаар сургалт, баталгаажуулалт, тестийн хэсэгт хуваан ашигласан нь загварын ерөнхий чадварыг үнэлэх боломжийг бүрдүүлнэ.
    \item \textbf{Ёс зүйн шинжилгээ:} Хүйсийн хазайлт, орчны хазайлт, арьс өнгийн хазайлт зэргийг тусгайлан судалсан нь судалгааны стандартыг дэмжинэ.
\end{itemize}

%-------------------------------------------------------------------------------
\section{Цаашдын судалгааны чиглэл}
%-------------------------------------------------------------------------------

Энэхүү судалгааны үр дүнд тулгуурлан дараах чиглэлүүдэд цаашдын судалгааг хийх нь зүйтэй:

\begin{enumerate}
    \item \textbf{Монгол хэлний тандалтын өгөгдлийн багц бэлтгэх:} Монгол хяналтын камерын бодит зодооны бичлэгийг цуглуулж, баяжуулсан өгөгдлийн багц бэлтгэх нь загварын гүйцэтгэлийг эрс нэмэгдүүлэх боломжтой. Монгол улсын хяналтын байгууллагуудтай хамтран ажиллах нь нэн чухал.

    \item \textbf{Загварын хэмжээ нэмэгдүүлэх:} Qwen2.5-14B загварт нарийн тохируулга хийж Qwen2.5-7B-тай харьцуулах --- нэмэлт гүйцэтгэл (7\%+) авах боломж байгаа эсэхийг тодорхойлох.

    \item \textbf{Бусад аномалийн ангилалд өргөтгөх:} Зодооны илрүүлэлтийг хулгай, осол, үймээн самуун зэрэг бусад аномалийн ангилалд өргөтгөн ерөнхий хэвийн бус зан үйл илрүүлэлтийн (general anomaly detection) систем бий болгох.

    \item \textbf{Шууд видео урсгал боловсруулалт:} Одоогийн тайлбар тайланд суурилсан аргаас видеог шууд боловсруулах (direct video дүгнэлт гаргалт) руу шилжүүлэх нь мэдээллийн алдагдлыг бууруулна.

    \item \textbf{Санал хүсэлтэд суурилсан идэвхтэй сургалт:} Операторын засварлалтыг ашиглан загварыг тасралтгүй сайжруулах (online learning) систем хэрэгжүүлэх.
\end{enumerate}
